\clearpage
\SourceLeaf{203}
\fancyhead[L]{\small Allgemeine Lehrsätze über Anziehungs- und Abstossungs-Kräfte}
\renewcommand{\Article}[1]{\par\medskip\Needspace{4\baselineskip}\pdfbookmark[1]{Art. #1}{lehrsaetze-art-#1}\centerline{\textbf{#1.}}\nobreak\smallskip}
\pdfbookmark[0]{Allgemeine Lehrsätze über Anziehungs- und Abstossungs-Kräfte}{lehrsaetze-body}
\begin{center}
ALLGEMEINE LEHRSÄTZE\\
IN BEZIEHUNG AUF DIE IM VERKEHRTEN VERHÄLTNISSE\\
DES QUADRATS DER ENTFERNUNG WIRKENDEN\\
ANZIEHUNGS- UND ABSTOSSUNGS-KRÄFTE.
\end{center}
\Article{1}
Die Natur bietet uns mancherlei Erscheinungen dar, welche wir durch die Annahme von Kräften erklären, die von den kleinsten Theilen der Substanzen auf einander ausgeübt werden, und den Quadraten der gegenseitigen Entfernungen umgekehrt proportional sind.

Vor allen gehört hieher die allgemeine Gravitation Vermöge derselben übt jedes ponderable Molecül $\mu$ auf ein anderes $\mu'$ eine bewegende Kraft aus, welche, wenn man die Entfernung $=r$ setzt, durch $\frac{\mu\mu'}{rr}$ ausgedrückt wird, und eine Annäherung in der Richtung der verbindenden geraden Linie hervorzubringen strebt.

Wenn man zur Erklärung der magnetischen Erscheinungen zwei magnetische Flüssigkeiten annimmt, wovon die eine als positive Grösse, die andere als negative betrachtet wird, so üben zwei derartige Elemente $\mu$, $\mu'$ gleichfalls eine bewegende Kraft auf einander aus, welche durch $\frac{\mu\mu'}{rr}$ gemessen wird, und in der verbindenden geraden Linie wirkt, aber als Abstossung, wenn $\mu$, $\mu'$ gleichartig, als Anziehung, wenn sie ungleichartig sind.

Ganz ähnliches gilt von der gegenseitigen Wirkung der Theile der elektrischen Flüssigkeiten auf einander.

\SourceLeaf{204}
Das linearische Element $\mathrm{d}s$ eines galvanischen Stroms übt auf ein Element des magnetischen Fluidums $\mu$ (wenn wir letzteres zulassen) ebenfalls eine bewegende Kraft aus, die dem Quadrate der Entfernung $r$ umgekehrt proportional ist. aber hier tritt zugleich der ganz abweichende Umstand ein, dass die Richtung der Kraft nicht in der verbindenden geraden Linie, sondern senkrecht gegen die durch $\mu$ und die Richtung von $\mathrm{d}s$ gelegte Ebene ist, und dass ausserdem die Stärke der Kraft nicht von der Entfernung allein, sondern zugleich von dem Winkel abhängt, welchen $r$ mit der Richtung von $\mathrm{d}s$ macht. Nennt man diesen Winkel $\theta$, so ist $\frac{\sin\theta.\mu\mathrm{d}s}{rr}$ das Maass der bewegenden Kraft, welche $\mathrm{d}s$ auf $\mu$ ausübt, und eben so gross ist die von $\mu$ auf das Stromelement $\mathrm{d}s$ oder dessen ponderabeln Träger ausgeübte Kraft, deren Richtung der erstern entgegengesetzt parallel ist.

Wenn man mit \textsc{Ampère} annimmt, dass zwei Elemente von galvanischen Strömen $\mathrm{d}s$, $\mathrm{d}s'$ in der sie verbindenden geraden Linie anziehend oder abstossend auf einander wirken, so nöthigen uns die Erscheinungen, diese Kraft gleichfalls dem Quadrate der Entfernung umgekehrt proportional zu setzen, zugleich aber erfordern jene eine etwas verwickeltere Abhängigkeit von der Richtung der Stromelemente.

Wir werden uns in dieser Abhandlung auf die drei ersten Fälle oder auf solche Kräfte einschränken, die sich in der Richtung der geraden Linie zwischen dem Elemente, welches wirkt, und demjenigen, auf welches gewirkt wird, äussern, und schlechthin dem Quadrate der Entfernung umgekehrt proportional sind, obwohl mehrere Lehrsätze mit geringer Veränderung auch bei den andern Fällen ihre Anwendung finden, deren ausführliche Entwickelung einer andern Abhandlung vorbehalten bleiben muss.

\Article{2}
Wir bezeichnen mit $a$, $b$, $c$ die rechtwinkligen Coordinaten eines materiellen Punktes, von welchem aus eine abstossende oder anziehende Kraft wirkt; die beschleunigende Kraft selbst in einem unbestimmten Punkte $O$, dessen Coordinaten $x$, $y$, $z$ sind, mit
\[
\frac{\mu}{(a-x)^2+(b-y)^2+(c-z)^2}=\frac{\mu}{rr}
\]
wo also $\mu$ für den ersten Fall des vorhergehenden Artikels die im erstern Punkte
\SourceLeaf{205}
befindliche ponderable Materie, im zweiten und dritten das Quantum magnetischen oder elektrischen Fluidums ausdrückt. Wird diese Kraft parallel mit den drei Coordinatenaxen zerlegt, so entstehen daraus die Componenten
\[
\frac{\varepsilon\mu(a-x)}{r^3},\qquad \frac{\varepsilon\mu(b-y)}{r^3},\qquad \frac{\varepsilon\mu(c-z)}{r^3}
\]
wo $\varepsilon=+1$ oder $=-1$ sein soll, jenachdem die Kraft anziehend oder abstossend wirkt, was sich nach der Beschaffenheit des Wirkenden und des die Wirkung Empfangenden von selbst entscheidet. Diese Componenten stellen sich dar als die partiellen Differentialquotienten
\[
\frac{\mathrm{d}\frac{\varepsilon\mu}{r}}{\mathrm{d}x},\qquad \frac{\mathrm{d}\frac{\varepsilon\mu}{r}}{\mathrm{d}y},\qquad \frac{\mathrm{d}\frac{\varepsilon\mu}{r}}{\mathrm{d}z}
\]
Wirken also auf denselben Punkt $O$ mehrere Agentien $\mu^0$, $\mu'$, $\mu''$ u. s. f. aus den Entfernungen $r^0$, $r'$, $r''$ u. s. f., und setzt man
\[
\frac{\mu^0}{r^0}+\frac{\mu'}{r'}+\frac{\mu''}{r''}+\text{ u. s. f.}=\Sigma\frac{\mu}{r}=V
\]
so werden die Componenten der ganzen in $O$ wirkenden Kraft durch
\[
\frac{\varepsilon\mathrm{d}V}{\mathrm{d}x},\qquad \frac{\varepsilon\mathrm{d}V}{\mathrm{d}y},\qquad \frac{\varepsilon\mathrm{d}V}{\mathrm{d}z}
\]
dargestellt.

Wenn die Agentien nicht aus discreten Punkten wirken, sondern eine Linie, eine Fläche oder einen körperlichen Raum stetig erfüllen, so tritt an die Stelle der Summation $\Sigma$ eine einfache, doppelte oder dreifache Integration. Der letzte Fall ist an sich allein der Fall der Natur: allein da man oft dafür, unter gewissen Einschränkungen, fingirte in Punkte concentrirte, oder auf Linien oder Flächen stetig vertheilte Agentien substituiren kann, so werden wir jene Fälle mit in unsre Untersuchung ziehen, wobei es unanstössig sein wird, von \textit{Massen}, die auf eine Fläche oder Linie vertheilt, oder in einen Punkt concentrirt sind, zu reden, insofern der Ausdruck Masse hier nichts weiter bedeutet, als dasjenige, wovon Anziehungs- oder Abstossungs-Kräfte ausgehend gedacht werden.

\Article{3}
Indem wir also, für jeden Punkt im Raume, mit $x$, $y$, $z$ dessen rechtwinklige Coordinaten, und mit $V$ das Aggregat aller wirkenden Massentheilchen, je\SourceLeaf{206}des mit seiner Entfernung von jenem Punkte dividirt, bezeichnen, wobei nach den jedesmaligen Bedingungen der Untersuchung negative Massentheilchen entweder ausgeschlossen oder als zulässig betrachtet werden mögen, wird $V$ eine Function von $x$, $y$, $z$, und die Erforschung der Eigenthümlichkeiten dieser Function der Schlüssel zur Theorie der Anziehungs- oder Abstossungskräfte selbst sein. Zur bequemern Handhabung der dazu dienenden Untersuchungen werden wir uns erlauben, dieses $V$ mit einer besondern Benennung zu belegen, und die Grösse das \textit{Potential} der Massen, worauf sie sich bezieht, nennen. Für unsre gegenwärtige Untersuchung reicht diese beschränktere Begriffsbestimmung hin: im weitern Sinn könnte man sowohl für Betrachtung anderer Anziehungsgesetze, als im umgekehrten Verhältniss des Quadrates der Entfernung, als auch für den vierten im Art. 1 erwähnten Fall, unter Potential die Function von $x$, $y$, $z$ verstehen, deren partielle Differentialquotienten die Componenten der erzeugten Kraft vorstellen.

Bezeichnen wir die ganze in dem Punkte $x$, $y$, $z$ Statt findende Kraft mit $p$, und die Winkel, welche ihre Richtung mit den drei Coordinatenaxen macht, mit $\alpha$, $\beta$, $\gamma$, so sind die drei Componenten
\[
p\cos\alpha=\varepsilon\frac{\mathrm{d}V}{\mathrm{d}x},\qquad p\cos\beta=\varepsilon\frac{\mathrm{d}V}{\mathrm{d}y},\qquad p\cos\gamma=\varepsilon\frac{\mathrm{d}V}{\mathrm{d}z}
\]
und
\[
p=\BareRoot{((\frac{\mathrm{d}V}{\mathrm{d}x})^2+(\frac{\mathrm{d}V}{\mathrm{d}y})^2+(\frac{\mathrm{d}V}{\mathrm{d}z})^2)}
\]
\Article{4}
Ist $\mathrm{d}s$ das Element einer beliebigen geraden oder krummen Linie, so sind $\frac{\mathrm{d}x}{\mathrm{d}s}$, $\frac{\mathrm{d}y}{\mathrm{d}s}$, $\frac{\mathrm{d}z}{\mathrm{d}s}$ die Cosinus der Winkel, welche jenes Element mit den Coordinatenaxen macht; bezeichnet also $\theta$ den Winkel zwischen der Richtung des Elements und der Richtung, welche die resultirende Kraft daselbst hat, so ist
\[
\cos\theta=\frac{\mathrm{d}x}{\mathrm{d}s}.\cos\alpha+\frac{\mathrm{d}y}{\mathrm{d}s}.\cos\beta+\frac{\mathrm{d}z}{\mathrm{d}s}.\cos\gamma
\]
Die auf die Richtung von $\mathrm{d}s$ projicirte Kraft wird folglich
\[
p\cos\theta=\varepsilon(\frac{\mathrm{d}V}{\mathrm{d}x}.\frac{\mathrm{d}x}{\mathrm{d}s}+\frac{\mathrm{d}V}{\mathrm{d}y}.\frac{\mathrm{d}y}{\mathrm{d}s}+\frac{\mathrm{d}V}{\mathrm{d}z}.\frac{\mathrm{d}z}{\mathrm{d}s})=\frac{\varepsilon\mathrm{d}V}{\mathrm{d}s}
\]
Legen wir durch alle Punkte, in welchen das Potential $V$ einen constanten Werth hat, eine Fläche, so wird solche, allgemein zu reden, die Theile des
\SourceLeaf{207}
Raums, wo $V$ kleiner ist, von denen scheiden, wo $V$ grösser ist als jener Werth. Liegt die Linie $s$ in dieser Fläche, oder tangirt sie wenigstens dieselbe mit dem Element $\mathrm{d}s$, so ist $\frac{\mathrm{d}V}{\mathrm{d}s}=0$. Falls also nicht an diesem Platze die Bestandtheile der ganzen Kraft einander destruiren, oder $p=0$ wird, in welchem Falle von einer Richtung der Kraft nicht mehr die Rede sein kann, muss nothwendig $\cos\theta=0$ sein, woraus wir schliessen, dass die Richtung der resultirenden Kraft in jedem Punkte einer solchen Fläche gegen diese selbst normal ist, und zwar nach derjenigen Seite des Raumes zu, wo die grössern Werthe von $V$ angrenzen, wenn $\varepsilon=+1$ ist; nach der entgegengesetzten, wenn $\varepsilon=-1$ ist. Wir nennen eine solche Fläche eine \textit{Gleichgewichtsfläche}. Da durch jeden Punkt eine solche Fläche gelegt werden kann, so wird die Linie $s$, falls sie nicht ganz in Einer Gleichgewichtsfläche liegt, in jedem ihrer Punkte eine andere treffen. Durchschneidet $s$ alle Gleichgewichtsflächen unter rechten Winkeln, so stellt eine Tangente an jener Linie überall die Richtung der Kraft, und $\frac{\mathrm{d}V}{\mathrm{d}s}$ ihre Stärke dar.

Das Integral $\int p\cos\theta.\mathrm{d}s$, durch ein beliebiges Stück der Linie $s$ ausgedehnt, wird offenbar $=\varepsilon(V'-V^0)$, wenn $V^0$, $V'$ die Werthe des Potentials für den Anfangs- und Endpunkt bedeuten. Ist also $s$ eine geschlossene Linie, so wird jenes Integral, durch die ganze Linie erstreckt, $=0$ werden.

\Article{5}
Es ist von selbst klar, dass das Potential in jedem Punkte des Raumes, der \textit{ausserhalb} aller anziehenden oder abstossenden Theilchen liegt, einen assignabeln Werth erhalten muss; dasselbe gilt aber auch von dessen Differentialquotienten, sowohl erster als höherer Ordnung, da diese in jener Voraussetzung gleichfalls die Form von Summen assignabler Theile oder von Integralen solcher Differentiale annehmen, in denen die Coëfficienten durchaus assignable Werthe haben. So wird
\[
\begin{array}{ll}
\displaystyle\frac{\mathrm{d}V}{\mathrm{d}x}=\Sigma\frac{(a-x)\mu}{r^3},&\displaystyle\frac{\mathrm{d}\mathrm{d}V}{\mathrm{d}x^2}=\Sigma(\frac{3(a-x)^2}{r^5}-\frac{1}{r^3})\mu\\[6pt]
\displaystyle\frac{\mathrm{d}V}{\mathrm{d}y}=\Sigma\frac{(b-y)\mu}{r^3},&\displaystyle\frac{\mathrm{d}\mathrm{d}V}{\mathrm{d}y^2}=\Sigma(\frac{3(b-y)^2}{r^5}-\frac{1}{r^3})\mu\\[6pt]
\displaystyle\frac{\mathrm{d}V}{\mathrm{d}z}=\Sigma\frac{(c-z)\mu}{r^3},&\displaystyle\frac{\mathrm{d}\mathrm{d}V}{\mathrm{d}z^2}=\Sigma(\frac{3(c-z)^2}{r^5}-\frac{1}{r^3})\mu
\end{array}
\]
\SourceLeaf{208}
Die bekannte Gleichung
\[
\frac{\mathrm{d}\mathrm{d}V}{\mathrm{d}x^2}+\frac{\mathrm{d}\mathrm{d}V}{\mathrm{d}y^2}+\frac{\mathrm{d}\mathrm{d}V}{\mathrm{d}z^2}=0
\]
gilt also für alle Punkte des Raumes, die ausserhalb der wirkenden Massen liegen.

\Article{6}
Unter den verschiedenen Fällen, wo der Werth des Potentials $V$ oder seiner Differentialquotienten für einen nicht ausserhalb der wirkenden Massen liegenden Punkt in Frage kommt, wollen wir zuerst den Fall der Natur betrachten, wo die Massen einen bestimmten körperlichen Raum mit gleichförmiger oder ungleichförmiger, aber überall endlicher Dichtigkeit ausfüllen.

Es sei $t$ der ganze Raum, welcher Masse enthält; $\mathrm{d}t$ ein unendlich kleines Element desselben, welchem die Coordinaten $a$, $b$, $c$ und das Massenelement $k\mathrm{d}t$ entsprechen; ferner sei $V$ das Potential in dem Punkte $O$, dessen Coordinaten $x$, $y$, $z$, also die Entfernung von jenem Element
\[
\BareRoot{((a-x)^2+(b-y)^2+(c-z)^2)}=r
\]
Es wird folglich
\[
V=\int\frac{k\mathrm{d}t}{r}
\]
durch den ganzen Raum $t$ ausgedehnt, was eine dreifache Integration implicirt. Man sieht leicht, dass eine wahre Integration stattnehmig ist, auch wenn $O$ innerhalb des Raumes sich befindet, obgleich dann $\frac{1}{r}$ für die unendlich nahe bei $O$ liegenden Elemente unendlich gross wird. Denn wenn man anstatt $a$, $b$, $c$ Polarcoordinaten einführt, indem man
\[
a=x+r\cos u,\qquad b=y+r\sin u\cos\lambda,\qquad c=z+r\sin u\sin\lambda
\]
setzt, so wird $\mathrm{d}t=rr\sin u.\mathrm{d}u.\mathrm{d}\lambda.\mathrm{d}r$, mithin
\[
V=\int\!\int\!\int kr\sin u.\mathrm{d}u.\mathrm{d}\lambda.\mathrm{d}r
\]
wo die Integration in Beziehung auf $r$ von $r=0$ bis zu dem an der Grenze von $t$ Statt findenden Werthe, von $\lambda=0$ bis $\lambda=2\pi$, und von $u=0$ bis $u=\pi$ ausgedehnt werden muss. Es wird also nothwendig $V$ einen bestimmten endlichen Werth erhalten.

\SourceLeaf{209}
Man sieht ferner leicht ein, dass man auch hier
\[
\frac{\mathrm{d}V}{\mathrm{d}x}=\int k\mathrm{d}t.\frac{\mathrm{d}\frac{1}{r}}{\mathrm{d}x}=\int\frac{k(a-x)\mathrm{d}t}{r^3}=X
\]
setzen darf. Die Befugniss dazu beruhet darauf, dass auch dieser Ausdruck, welcher unter Anwendung von Polarcoordinaten in
\[
\int\!\int\!\int k\cos u.\sin u.\mathrm{d}u.\mathrm{d}\lambda.\mathrm{d}r
\]
übergeht, einer wahren Integration fähig ist, also $X$ einen bestimmten endlichen Werth erhält, der sich nach der Stetigkeit ändert, weil alle in unendlicher Nähe bei $O$ liegenden Elemente nur einen unendlich kleinen Beitrag dazu geben. Aus ähnlichen Gründen darf man auch
\[
\begin{array}{l}
\displaystyle\frac{\mathrm{d}V}{\mathrm{d}y}=\int\frac{k(b-y)\mathrm{d}t}{r^3}=Y\\[6pt]
\displaystyle\frac{\mathrm{d}V}{\mathrm{d}z}=\int\frac{k(c-z)\mathrm{d}t}{r^3}=Z
\end{array}
\]
setzen, und diese Grössen erhalten daher, eben so wie $V$, innerhalb $t$ bestimmte nach der Stetigkeit sich ändernde Werthe. Dasselbe wird auch noch auf der Grenze von $t$ gelten.

\Article{7}
Was nun aber die Differentialquotienten höherer Ordnungen betrifft, so muss für Punkte innerhalb $t$ ein anderes Verfahren eintreten, da es z. B. nicht verstattet ist, $\frac{\mathrm{d}X}{\mathrm{d}x}$ in
\[
\int k\mathrm{d}t.\frac{\mathrm{d}\frac{a-x}{r^3}}{\mathrm{d}x}
\]
d. i. in
\[
\int k(\frac{3(a-x)^2-rr}{r^5})\mathrm{d}t
\]
umzuformen, indem dieser Ausdruck genau betrachtet nur ein Zeichen ohne bestimmte klare Bedeutung sein würde. Denn in der That, da sich innerhalb jedes auch noch so kleinen Theils von $t$, welcher den Punkt einschliesst, Theile nachweisen lassen, über welche ausgedehnt dieses Integral jeden vorgegebenen Werth, er sei positiv oder negativ, überschreitet, so fehlt hier die wesentliche
\SourceLeaf{210}
Bedingung, unter welcher allein dem ganzen Integrale eine klare Bedeutung beigelegt werden kann, nemlich die Anwendbarkeit der Exhaustionsmethode.

\Article{8}
Ehe wir diese Untersuchung in ihrer Allgemeinheit vornehmen, wird es zur Fixirung der Vorstellungen nützlich sein, einen sehr einfachen speciellen Fall zu betrachten.

Es sei $t$ eine Kugel, deren Halbmesser $=R$ ist, und deren Mittelpunkt mit dem Anfangspunkte der Coordinaten zusammenfällt: die Dichtigkeit der die Kugel erfüllenden Masse sei constant $=k$, und den Abstand des Punktes $O$ vom Mittelpunkte bezeichnen wir mit $\rho=\BareRoot{(xx+yy+zz)}$. Bekanntlich hat das Potential zwei verschiedene Ausdrücke, je nachdem $O$ innerhalb der Kugel, oder ausserhalb liegt. Im erstern Fall ist nemlich
\[
V=2\pi kRR-\tfrac{2}{3}\pi k\rho\rho=2\pi kRR-\tfrac{2}{3}\pi k(xx+yy+zz)
\]
im zweiten hingegen
\[
V=\frac{4\pi kR^3}{3\rho}
\]
Auf der Oberfläche der Kugel geben beide Ausdrücke einerlei Werth $\tfrac{4}{3}\pi kRR$, und das Potential ändert sich daher im ganzen Raume nach der Stetigkeit.

Für die Differentialquotienten erhalten wir, im innern Raume
\[
\begin{array}{l}
\displaystyle\frac{\mathrm{d}V}{\mathrm{d}x}=X=-\tfrac{4}{3}\pi kx\\[6pt]
\displaystyle\frac{\mathrm{d}V}{\mathrm{d}y}=Y=-\tfrac{4}{3}\pi ky\\[6pt]
\displaystyle\frac{\mathrm{d}V}{\mathrm{d}z}=Z=-\tfrac{4}{3}\pi kz
\end{array}
\]
im äussern Raume hingegen
\[
\begin{array}{l}
\displaystyle X=-\frac{4\pi kR^3x}{3\rho^3}\\[6pt]
\displaystyle Y=-\frac{4\pi kR^3y}{3\rho^3}\\[6pt]
\displaystyle Z=-\frac{4\pi kR^3z}{3\rho^3}
\end{array}
\]
Auch hier geben auf der Oberfläche die letztern Formeln dieselben Werthe wie die erstern, daher auch $X$, $Y$, $Z$ im ganzen Raume nach der Stetigkeit sich ändern.

\SourceLeaf{211}
Anders verhält es sich aber mit den Differentialquotienten dieser Grössen. Im innern Raume haben wir
\[
\frac{\mathrm{d}X}{\mathrm{d}x}=-\tfrac{4}{3}\pi k,\qquad \frac{\mathrm{d}Y}{\mathrm{d}y}=-\tfrac{4}{3}\pi k,\qquad \frac{\mathrm{d}Z}{\mathrm{d}z}=-\tfrac{4}{3}\pi k
\]
im äussern Raume hingegen
\[
\begin{array}{l}
\displaystyle\frac{\mathrm{d}X}{\mathrm{d}x}=\frac{4\pi kR^3(3xx-\rho\rho)}{3\rho^5}\\[6pt]
\displaystyle\frac{\mathrm{d}Y}{\mathrm{d}y}=\frac{4\pi kR^3(3yy-\rho\rho)}{3\rho^5}\\[6pt]
\displaystyle\frac{\mathrm{d}Z}{\mathrm{d}z}=\frac{4\pi kR^3(3zz-\rho\rho)}{3\rho^5}
\end{array}
\]
Auf der Oberfläche fallen diese Werthe nicht mit jenen zusammen, sondern sind beziehungsweise
\[
\frac{4\pi kxx}{RR},\qquad \frac{4\pi kyy}{RR},\qquad \frac{4\pi kzz}{RR}
\]
grösser. Es ändern sich daher jene Differentialquotienten, nach der Stetigkeit zwar im ganzen innern und im ganzen äussern Raume, aber sprungsweise beim Übergange aus dem einen in den andern, und in der Scheidungsfläche selbst muss man ihnen doppelte Werthe beilegen, je nachdem $\mathrm{d}x$, $\mathrm{d}y$, $\mathrm{d}z$ als positiv oder als negativ betrachtet werden.

Ähnliches findet bei den sechs übrigen Differentialquotienten
\[
\frac{\mathrm{d}X}{\mathrm{d}y},\qquad \frac{\mathrm{d}X}{\mathrm{d}z},\qquad \frac{\mathrm{d}Y}{\mathrm{d}x},\qquad \frac{\mathrm{d}Y}{\mathrm{d}z},\qquad \frac{\mathrm{d}Z}{\mathrm{d}x},\qquad \frac{\mathrm{d}Z}{\mathrm{d}y}
\]
Statt, die im Innern der Kugel sämmtlich $=0$ werden, und beim Durchgange durch die Kugelfläche sprungsweise die Änderungen
\[
\frac{4\pi kxy}{RR},\qquad \frac{4\pi kxz}{RR}\text{ u. s. f.}
\]
erleiden.

Das Aggregat $\frac{\mathrm{d}X}{\mathrm{d}x}+\frac{\mathrm{d}Y}{\mathrm{d}y}+\frac{\mathrm{d}Z}{\mathrm{d}z}$ oder $\frac{\mathrm{d}\mathrm{d}V}{\mathrm{d}x^2}+\frac{\mathrm{d}\mathrm{d}V}{\mathrm{d}y^2}+\frac{\mathrm{d}\mathrm{d}V}{\mathrm{d}z^2}$ wird im Innern der Kugel $=-4\pi k$, im äussern Raume $=0$. Auf der Oberfläche selbst verliert es aber seine einfache Bedeutung: präcis zu reden, kann man nur sagen, dass es ein Aggregat von drei Theilen ist, deren jeder zwei verschiedene Werthe hat, und so gibt es eigentlich acht Combinationen, unter denen eine mit dem auf der innern Seite, eine andere mit dem auf der äussern Seite geltenden Werthe übereinstimmt, während die sechs übrigen ohne alle Bedeutung bleiben. Der Ana%
